\documentclass[11pt,a4paper]{article}

% --- CẤU HÌNH NGÔN NGỮ VÀ MÃ HÓA ---
\usepackage{iftex}
\ifxetex
  \usepackage{fontspec}
  \setmainfont{Times New Roman}
\else
  \usepackage[utf8]{vietnam}
\fi

% --- CẤU HÌNH LỀ TRANG & ĐỊNH DẠNG ---
\usepackage[top=2.0cm, bottom=2.0cm, left=2.5cm, right=2.0cm]{geometry}
\usepackage{amsmath, amssymb, amsfonts}
\usepackage{booktabs, multirow, array}
\usepackage{graphicx}
\usepackage{xcolor}
\usepackage{hyperref}
\usepackage{cite}
\usepackage{caption}
\usepackage{subcaption}

% Cấu hình liên kết tài liệu
\hypersetup{
    colorlinks=true,
    linkcolor=blue!70!black,
    citecolor=blue!70!black,
    urlcolor=blue!70!black
}

% Cấu hình dãn dòng
\renewcommand{\baselinestretch}{1.2}
\setlength{\parskip}{0.5em}
\setlength{\parindent}{1.5em}

% --- THÔNG TIN BÁO CÁO ---
\title{\textbf{\Large BÁO CÁO KHOA HỌC: PHƯƠNG PHÁP LUẬN VÀ THIẾT KẾ KIẾN TRÚC\\KHAI THÁC CẶP ẢNH NỀN TĨNH VÀ PHƯƠNG TIỆN TRONG THỊ GIÁC GIAO THÔNG VỚI DINO}}
\author{\textbf{Nhóm Nghiên Cứu Thị Giác Máy Tính Đô Thị}}
\date{\today}

\begin{document}

\maketitle

\begin{abstract}
Trong hệ thống camera giám sát giao thông thông minh (CCTV), các thiết bị ghi hình cố định cung cấp cặp tín hiệu quang học tự nhiên có giá trị: ảnh hiện trường ghi nhận phương tiện di chuyển và ảnh nền tĩnh (background) sạch bóng xe được ước lượng qua chuỗi thời gian. Báo cáo này trình bày cơ sở lý thuyết, mô hình hóa toán học và phương pháp luận chuyên sâu của bảy hướng nghiên cứu học tự giám sát (Self-Supervised Learning) nhằm khai thác triệt để cặp tín hiệu trên. Hệ sinh thái nghiên cứu bao gồm: (1) Tiền huấn luyện Vision Transformer với cơ chế che có định hướng tiền cảnh (Foreground-Aware Masking); (2) Mạng phân rã bối cảnh tự giám sát (Scene Decomposition Network) dựa trên mô hình hòa trộn Alpha Compositing và xóa xe tự động (Inpainting); (3) Mở rộng biểu diễn Patch Embedding bốn kênh (RGB + $\Delta$) phục vụ bài toán ước lượng lưu lượng phương tiện trong điều kiện ít mẫu (Few-shot learning); (5) Dự đoán mật độ không-thời gian và phân loại cấp độ dịch vụ giao thông theo chuẩn Highway Capacity Manual (HCM LoS) dựa trên mạng hợp nhất DINO ViT, CNN và Bi-GRU; (6) Nhận dạng lại phương tiện liên camera theo hành lang (Corridor-based Vehicle Re-ID) không cần Object Detector kết hợp cửa sổ khả thi không-thời gian nhằm ước tính thời gian hành trình và tốc độ di chuyển; (7) Hiểu cảnh giao thông từ vựng mở (Open-Vocabulary Scene Understanding) kết hợp đề xuất vùng vật lý $\Delta$ và không gian ngữ nghĩa CLIP; và (8) Đánh giá đa thuộc tính điều kiện mặt đường tự giám sát từ chuỗi ảnh nền 24 giờ. Nghiên cứu loại bỏ các tiếp cận bất thường không có mỏ neo vật lý để tập trung vào các bài toán đo lường động học giao thông có tính ứng dụng thực tiễn cao nhất.
\end{abstract}

\hrule
\vspace{1em}

\section{ĐẶT VẤN ĐỀ VÀ MÔ HÌNH HÓA TÍN HIỆU QUANG HỌC}

Trong bài toán giám sát giao thông đô thị qua camera tĩnh, mỗi khung hình tại thời điểm $t$ được xem xét dưới hai trạng thái quan sát:
\begin{enumerate}
    \item \textbf{Khung hình hiện trường ($I_{\text{origin}} \in \mathbb{R}^{H \times W \times 3}$):} Ghi nhận các phương tiện giao thông (xe máy, ô tô, xe buýt, xe tải) chuyển động trên bề mặt tuyến đường, kèm theo các yếu tố biến đổi về mật độ, điều kiện chiếu sáng và mức độ che khuất lẫn nhau.
    \item \textbf{Ảnh nền tĩnh ($I_{\text{bg}} \in \mathbb{R}^{H \times W \times 3}$):} Thể hiện cấu trúc mặt đường khi không có phương tiện, được ước lượng thông qua thuật toán lọc trung vị theo thời gian (Temporal Median Filtering) kết hợp chuẩn hóa cường độ trong không gian màu HSV theo từng chu kỳ quan sát.
\end{enumerate}

Mối quan hệ không gian giữa hai khung hình xác lập trường sai khác quang học:
\begin{equation}
    \Delta(u, v) = \|I_{\text{origin}}(u, v) - I_{\text{bg}}(u, v)\|
\end{equation}

Thay vì chỉ áp dụng $\Delta(u, v)$ trong các thuật toán phân đoạn ngưỡng cổ điển vốn nhạy cảm với biến thiên thời tiết và bóng đổ, hệ sinh thái DINO Suite tích hợp tín hiệu sai khác vật lý này vào không gian biểu diễn sâu của Vision Transformer (ViT), mở ra bảy hướng tiếp cận bổ trợ lẫn nhau trong giao thông thông minh.

\section{CƠ SỞ PHƯƠNG PHÁP LUẬN VỀ HỌC TỰ GIÁM SÁT DINO}

\subsection{Cơ Chế Chưng Cất Tự Thân và Động Học Ổn Định Biểu Diễn}
Kiến trúc DINO (Caron et al., ICCV 2021) dựa trên nguyên lý tự chưng cất tri thức (self-distillation) giữa hai mạng nơ-ron có cấu trúc tham số tương đồng: mạng Student $g_{\theta_s}$ và mạng Teacher $g_{\theta_t}$. Mỗi mạng bao gồm một bộ trích xuất đặc trưng Vision Transformer (ViT) kết hợp đầu chiếu phi tuyến đa tầng (MLP Projection Head) nhằm ánh xạ vector đặc trưng $[CLS]$ sang không gian phân phối xác suất $K$ chiều ($K \in \{4096, 65536\}$).

Hệ thống tối ưu hóa vận hành theo quy tắc cập nhật tham số bất đối xứng:
\begin{enumerate}
    \item \textbf{Tối ưu hóa Student qua lan truyền ngược:} Tham số $\theta_s$ được cập nhật trực tiếp bằng thuật toán AdamW dựa trên độ dốc của hàm mất mát entropy chéo giữa đầu ra Student và Teacher.
    \item \textbf{Cập nhật Teacher theo trung bình trượt hàm mũ (EMA):} Trọng số của Teacher không tiếp nhận trực tiếp gradient ($\nabla_{\theta_t} \mathcal{L} = 0$), mà tiến hóa theo trung bình trượt động lượng từ tham số Student:
    \begin{equation}
        \theta_t \leftarrow \lambda \theta_t + (1 - \lambda) \theta_s
    \end{equation}
    trong đó hệ số động lượng $\lambda \in [0.996, 1.0]$ được điều chỉnh tăng dần theo lịch trình cosin. Về mặt giải tích, cơ chế EMA đóng vai trò như bộ lọc thông thấp trong không gian tham số, giúp mạng Teacher cung cấp phân phối xác suất mục tiêu ổn định và có phương sai thấp qua các bước huấn luyện.
    \item \textbf{Cơ chế kiểm soát sự sụp đổ biểu diễn (Representation Collapse Prevention):} DINO kiểm soát trạng thái suy biến bằng việc kết hợp kỹ thuật định tâm (\textit{Centering}) và kỹ thuật làm nhọn phân phối (\textit{Sharpening}):
    \begin{equation}
        P_t(x)^{(k)} = \frac{\exp\left(\frac{g_{\theta_t}(x)^{(k)} - c^{(k)}}{\tau_t}\right)}{\sum_{j=1}^K \exp\left(\frac{g_{\theta_t}(x)^{(j)} - c^{(j)}}{\tau_t}\right)}, \quad P_s(x)^{(k)} = \frac{\exp\left(\frac{g_{\theta_s}(x)^{(k)}}{\tau_s}\right)}{\sum_{j=1}^K \exp\left(\frac{g_{\theta_s}(x)^{(j)}}{\tau_s}\right)}
    \end{equation}
    trong đó vector định tâm $c \in \mathbb{R}^K$ ước lượng kỳ vọng bậc một của phân phối đầu ra Teacher, và nhiệt độ $\tau_t < \tau_s$ tạo hiệu ứng làm nhọn phân phối xác suất.
\end{enumerate}

\subsection{Sự Phát Triển Kiến Trúc Từ DINOv1 Đến DINOv2}
DINOv2 (Oquab et al., Meta AI 2023) mở rộng khung lý thuyết DINO từ việc học biểu diễn toàn cục $[CLS]$ thành mô hình nền tảng thị giác (Vision Foundation Model) có khả năng sinh biểu diễn dày đặc (dense representations) ở cấp độ từng patch, với các cải tiến kiến trúc then chốt:
\begin{enumerate}
    \item \textbf{Cơ chế học đa nhiệm tích hợp Masked Image Modeling (MIM / iBOT):} DINOv2 kết hợp đồng thời việc tự chưng cất token toàn cục $[CLS]$ với bài toán phục hồi biểu diễn của các patch bị che ngẫu nhiên:
    \begin{equation}
        \mathcal{L}_{\text{total}} = \mathcal{L}_{\text{DINO}}([CLS]) + \lambda_{\text{patch}} \mathcal{L}_{\text{iBOT}}(\text{Patch Tokens})
    \end{equation}
    \item \textbf{Ổn định hóa truyền gradient trên mô hình quy mô lớn:} Áp dụng LayerScale, Pre-LayerNorm, và SwiGLU FFN.
    \item \textbf{Khử nhiễu đặc trưng cục bộ qua Register Tokens:} Bổ sung 4 register tokens vào chuỗi đầu vào để hấp thụ các tín hiệu kích hoạt ngoại lai trên nền trơn nhẵn (nhựa đường, bầu trời).
\end{enumerate}

\subsection{Chiến Lược Trích Xuất Đa Góc Nhìn (Multi-Crop)}
Chiến lược đa góc nhìn sinh ra tập hợp gồm $N_g = 2$ góc nhìn toàn cảnh ($224 \times 224$, tỷ lệ diện tích $[0.4, 1.0]$) nạp vào cả Teacher và Student, cùng $N_l = 4$ đến $8$ góc nhìn cục bộ ($96 \times 96$, tỷ lệ diện tích $[0.05, 0.4]$) chỉ nạp vào Student. 
Kỳ vọng toán học của độ giao thoa không gian giữa hai góc nhìn toàn cảnh ngẫu nhiên thỏa mãn $\mathbb{E}[\text{IoU}] \in [0.15, 0.60]$, ép mô hình vừa học tính bất biến quang học (khi $\text{IoU}$ lớn) vừa trừu tượng hóa ngữ cảnh toàn cục của đối tượng (khi $\text{IoU} \to 0$).

\section{HƯỚNG 1: BG-GUIDED DINO -- TIỀN HUẤN LUYỆN TỰ GIÁM SÁT LIÊN TỤC VỚI FOREGROUND-AWARE MASKING}

\subsection{Hạn Chế của Cơ Chế Che Ngẫu Nhiên Đồng Đều}
Trong ảnh camera giao thông đô thị, hơn $70\% - 80\%$ diện tích khung hình là nền đường tĩnh, vỉa hè và dải phân cách bê tông. Nếu áp dụng chiến lược che ngẫu nhiên đồng đều (Uniform Random Masking), phần lớn tài nguyên tính toán và gradient cập nhật của ViT bị lãng phí chỉ để tái tạo các mảng nhựa đường vô nghĩa.

\subsection{Cơ Sở Toán Học của Chiến Lược Che Foreground-Aware Masking (FAM)}
Ảnh được chia thành $N_{\text{patches}} = \left(\frac{H}{P}\right) \times \left(\frac{W}{P}\right)$ patch rời rạc ($P=16$). Trọng số sai khác trung bình của patch thứ $p$ được tính trong không gian màu CIE-LAB:
\begin{equation}
    w_p = \frac{1}{P^2} \sum_{(u, v) \in \text{Patch}_p} \Delta(u, v)
\end{equation}

Xác suất để patch $p$ được lựa chọn che được xác định theo mô hình nội suy tuyến tính:
\begin{equation}
    P_{\text{mask}}(p) = \alpha \cdot \frac{w_p}{\max_{q} w_q + \epsilon} + (1 - \alpha) \cdot \frac{1}{N_{\text{patches}}}
\end{equation}
với $\alpha = 0.75$. Thành phần thứ nhất phân bổ $75\%$ xác suất che tập trung vào các vùng phương tiện (buộc Student phải suy diễn ngữ nghĩa xe từ ngữ cảnh xung quanh), trong khi thành phần thứ hai duy trì $25\%$ phân phối đều để bảo đảm mạng không bỏ sót cấu trúc không gian nền đường.

Hàm mất mát chưng cất đa góc nhìn:
\begin{equation}
    \mathcal{L}_{\text{DINO}} = -\sum_{x \in \{V^g_{\text{masked}}, V^l\}} \sum_{x' \in \{V_1^g, V_2^g\}, x' \neq x} P_t(x') \log P_s(x)
\end{equation}

\section{HƯỚNG 2: MẠNG PHÂN RÃ BỐI CẢNH TỰ GIÁM SÁT (TRAFFIC-DECOMPOSE)}

\subsection{Mô Hình Hóa Alpha Compositing với Mỏ Neo Nền Tĩnh}
Khung cảnh giao thông được mô hình hóa theo phương trình hòa trộn quang học:
\begin{equation}
    I_{\text{origin}} = M_{\alpha} \odot I_{\text{fg}} + (1 - M_{\alpha}) \odot I_{\text{bg}}
\end{equation}
trong đó $I_{\text{bg}}$ là lớp nền đường sạch bóng xe, $I_{\text{fg}}$ là lớp chứa phương tiện cô lập, và $M_{\alpha} \in [0, 1]^{H \times W \times 1}$ là mặt nạ mờ trong suốt (Alpha Matte).

Để giải bài toán nghịch đảo thiếu ràng buộc này, nghiên cứu sử dụng ảnh nền tĩnh thật $I_{\text{bg\_real}}$ làm mỏ neo vật lý giám sát tự thân. Mạng nơ-ron sử dụng chung bộ trích xuất đặc trưng ViT và phân nhánh thành ba đầu giải mã (Background Head, Foreground Head, Alpha Head).

\subsection{Hàm Mất Mát Đa Ràng Buộc Vật Lý}
\begin{equation}
    \mathcal{L}_{\text{total}} = \lambda_{\text{rec}} \mathcal{L}_{\text{recon}} + \lambda_{\text{bg}} \mathcal{L}_{\text{bg}} + \lambda_{\text{sparse}} \mathcal{L}_{\text{sparsity}} + \lambda_{\text{tv}} \mathcal{L}_{\text{tv}}
\end{equation}
\begin{itemize}
    \item $\mathcal{L}_{\text{recon}} = \|I_{\text{origin}} - \hat{I}_{\text{origin}}\|_{1} + \big(1 - \text{SSIM}(I_{\text{origin}}, \hat{I}_{\text{origin}})\big)$: Đảm bảo tái tạo hoàn hảo ảnh gốc.
    \item $\mathcal{L}_{\text{bg}} = \|\hat{I}_{\text{bg}} - I_{\text{bg\_real}}\|_{1}$: Khóa chặt nhánh nền vào lòng đường thực tế, buộc nhánh $I_{\text{fg}}$ và $M_\alpha$ phải hấp thụ toàn bộ phương tiện chuyển động.
    \item $\mathcal{L}_{\text{sparsity}} = \frac{1}{HW} \sum_{u, v} M_{\alpha}(u, v)$: Ngăn chặn nghiệm suy biến $M_{\alpha} = 1$.
    \item $\mathcal{L}_{\text{tv}} = \text{TotalVariation}(M_{\alpha})$: Khử nhiễu đốm và làm mịn đường biên thân xe.
\end{itemize}
Mô hình cho phép tự động xóa sạch phương tiện (Unsupervised Road Inpainting) trực tiếp từ một khung hình đơn lẻ.

\section{HƯỚNG 3: MỞ RỘNG BIỂU DIỄN TIỀN CẢNH BỐN KÊNH TRONG ƯỚC LƯỢNG LƯU LƯỢNG PHƯƠNG TIỆN}

\subsection{Mở Rộng Tensor Đầu Vào và Khởi Tạo Thích Ứng Ấm}
Bài toán ước lượng lưu lượng phương tiện đối mặt với hiện tượng che khuất nghiêm trọng trong giờ cao điểm và điều kiện học ít mẫu (Few-shot learning: $5\%, 10\%, 20\%$ nhãn). Tensor đầu vào được mở rộng thành bốn kênh:
\begin{equation}
    X_{\text{4ch}} = [\text{R}, \text{G}, \text{B}, \Delta_{\text{norm}}] \in \mathbb{R}^{4 \times H \times W}
\end{equation}

Nhằm bảo toàn tri thức tiền huấn luyện của ViT, thuật toán Warm-Start Weight Adaptation được áp dụng cho ma trận trọng số Patch Embedding $W_{\text{4ch}} \in \mathbb{R}^{D \times 4 \times P \times P}$:
\begin{equation}
    W_{\text{4ch}}[:, 0:3, :, :] = W_{\text{pretrained}}, \quad W_{\text{4ch}}[:, 3, :, :] = \frac{1}{3} \sum_{c=0}^{2} W_{\text{pretrained}}[:, c, :, :]
\end{equation}
Trọng số kênh $\Delta$ nhận trung bình cộng ba kênh màu RGB, bảo đảm tại bước huấn luyện đầu tiên kênh sai khác đóng góp năng lượng đồng mức mà không gây sốc gradient cho mạng nơ-ron. Đánh giá được thực hiện qua giao thức phân chia camera phân tách độc lập (Spatial Disjoint Splitting) nhằm triệt tiêu hiện tượng rò rỉ dữ liệu bối cảnh.

\section{HƯỚNG 5: DỰ ĐOÁN MẬT ĐỘ KHÔNG-THỜI GIAN VÀ CẤP ĐỘ DỊCH VỤ GIAO THÔNG (SPATIO-TEMPORAL DENSITY \& HCM LoS)}

\subsection{Đặt Vấn Đề và Mỏ Neo Vật Lý Tự Thân}
Các mô hình thị giác truyền thống thường tiếp cận theo hướng đếm phương tiện rời rạc trên từng khung hình tĩnh. Trong dòng giao thông xe máy hỗn hợp đặc trưng tại các đô thị Đông Nam Á (hàng trăm xe máy di chuyển sát nhau), hiện tượng che khuất chồng chéo dẫn đến sai số lũy kế nghiêm trọng. Bên cạnh đó, các tiếp cận học tương phản thời gian (Temporal InfoNCE) ngẫu nhiên không có mỏ neo vật lý thường bị mất tính hội tụ do biến thiên giao thông đột ngột giữa các khung giờ.

Nghiên cứu chuyển đổi mục tiêu sang ước lượng \textbf{Tỷ lệ Chiếm dụng Lòng đường Liên tục} $\rho(t) \in [0, 1]$ và \textbf{Cấp độ Dịch vụ Giao thông} theo chuẩn Highway Capacity Manual (HCM LoS).
Trường sai khác quang học $\Delta_t$ cung cấp mỏ neo vật lý tự thân mà không cần gán nhãn thủ công:
\begin{equation}
    \rho_{\text{phys}}(t) = \frac{1}{H \times W} \sum_{u=1}^H \sum_{v=1}^W \mathbb{I}\big(\Delta_t(u, v) > \tau\big)
\end{equation}
trong đó $\mathbb{I}(\cdot)$ là hàm chỉ thị và $\tau$ là ngưỡng phân tách quang học thích ứng.

\subsection{Rời Rạc Hóa Cấp Độ Dịch Vụ và Đạo Hàm Xu Hướng}
Giá trị $\rho_{\text{phys}}(t)$ được phân loại thành 4 cấp độ dịch vụ chuẩn HCM LoS:
\begin{equation}
    \text{LoS}(t) = \begin{cases} 
    0 \quad (\text{Free-Flow: Lưu thông tự do, thông thoáng}), & \rho(t) < 0.15 \\ 
    1 \quad (\text{Moderate: Dòng xe ổn định, mật độ trung bình}), & 0.15 \le \rho(t) < 0.35 \\ 
    2 \quad (\text{Slow: Mật độ cao, dòng xe di chuyển chậm}), & 0.35 \le \rho(t) < 0.60 \\ 
    3 \quad (\text{Gridlock: Kẹt xe nghiêm trọng, tê liệt}), & \rho(t) \ge 0.60 
    \end{cases}
\end{equation}
Đạo hàm xu hướng biến thiên mật độ được mô hình hóa qua sai phân hữu hạn: $\delta(t) = \rho(t) - \rho(t-1) \in [-1, 1]$.

\subsection{Kiến Trúc Hợp Nhất Không - Thời Gian (ViT + CNN + Bi-GRU)}
Mô hình tiếp nhận chuỗi gồm $T$ khung hình liên tiếp. Tại mỗi bước thời gian $t$:
\begin{enumerate}
    \item Vector ngữ nghĩa sâu: $\mathbf{z}_{\text{sem}}^t = \text{DINO}(I_t)_{[CLS]} \in \mathbb{R}^D$.
    \item Vector hình thái tiền cảnh: $\mathbf{z}_{\text{spatial}}^t = \text{CNN}(\Delta_t) \in \mathbb{R}^{128}$.
    \item Hợp nhất không gian: $\mathbf{x}_t = \text{Linear}([\mathbf{z}_{\text{sem}}^t \,\|\, \mathbf{z}_{\text{spatial}}^t]) \in \mathbb{R}^{256}$.
    \item Mô hình hóa chuỗi qua mạng Bi-GRU 2 tầng: $\mathbf{h}_t = \text{BiGRU}(\mathbf{x}_t, \mathbf{h}_{t-1}) \in \mathbb{R}^{256}$.
\end{enumerate}

Hàm mất mát đa nhiệm tích hợp thành phần làm mịn thời gian:
\begin{equation}
\begin{split}
    \mathcal{L}_{\text{total}} = \; & \lambda_{\rho} \mathcal{L}_{\text{SmoothL1}}(\hat{\rho}, \rho_{\text{phys}}) + \lambda_{\text{LoS}} \mathcal{L}_{\text{CE}}(\hat{\mathbf{y}}_{\text{LoS}}, y_{\text{LoS}}) \\
    & + \lambda_{\text{trend}} \mathcal{L}_{\text{SmoothL1}}(\hat{\delta}, \delta_{\text{phys}}) + \lambda_{\text{smooth}} \frac{1}{T-1} \sum_{t=1}^{T-1} \|\hat{\rho}_{t+1} - \hat{\rho}_t\|_2^2
\end{split}
\end{equation}

\section{HƯỚNG 6: ĐỊNH DANH LẠI PHƯƠNG TIỆN LIÊN CAMERA THEO HÀNH LANG VÀ ƯỚC TÍNH THỜI GIAN DI CHUYỂN}

\subsection{Đặt Vấn Đề và Hạn Chế của So Khớp Re-ID Truyền Thống}
Các giải pháp Vehicle Re-ID truyền thống thực hiện so khớp Cosine tự do trên toàn bộ thư viện Gallery. Trong bối cảnh đô thị Việt Nam với hàng chục nghìn xe máy có ngoại quan gần như đồng nhất, việc so khớp hình ảnh đơn thuần gây ra tỷ lệ \textbf{dương tính giả (False Positives) khổng lồ}. Hơn nữa, việc phụ thuộc vào Object Detector (YOLO) làm tiêu hao chi phí tính toán phần cứng lớn.

\subsection{Trích Xuất Vùng Xe Tự Động (\texorpdfstring{$\Delta$}{Delta}-RoI) và Kiến Trúc Biểu Diễn Danh Tính}
Bản đồ $\Delta$ được xử lý hình thái học lấp đầy thân xe và phân ngưỡng Otsu để xác định vùng liên thông hợp lệ không cần Object Detector:
\begin{equation}
    \text{RoI}_i = \text{ConnectedComponents}\big(\text{MorphClose}(\Delta > \tau_{\text{Otsu}})\big)_i
\end{equation}
Đặc trưng danh tính được trích xuất bằng DINO ViT kết hợp tầng gom cụm Generalized-Mean (GeM) Pooling và cấu trúc cổ chai BNNeck, tạo vector biểu diễn $\mathbf{e}_i \in \mathbb{R}^{256}$ chuẩn hóa $L_2$ ($\|\mathbf{e}_i\|_2 = 1$).

\subsection{Ràng Buộc Khả Thi Không - Thời Gian Theo Hành Lang (Spatio-Temporal Windowing)}
Xét hai camera liên tiếp $c_1$ và $c_2$ dọc theo cùng một hành lang giao thông với cự ly thực địa $d(c_1, c_2)$ (tính bằng km).
Dải vận tốc vật lý thực tế của phương tiện được giới hạn trong khoảng $[v_{\text{min}}, v_{\text{max}}] = [10, 60] \text{ km/h}$. Cửa sổ thời gian di chuyển hợp lệ bắt buộc thỏa mãn:
\begin{equation}
    \Delta t = t_2 - t_1 \in \left[ \Delta t_{\text{min}}, \Delta t_{\text{max}} \right] = \left[ \frac{d(c_1, c_2)}{v_{\text{max}}}, \ \frac{d(c_1, c_2)}{v_{\text{min}}} \right]
\end{equation}

Mọi ứng viên $g_j$ ngoài khoảng thời gian này bị triệt tiêu ngay lập tức khỏi phép so khớp:
\begin{equation}
    \text{Sim}(q_i, g_j) = \begin{cases} \mathbf{e}_{q_i} \cdot \mathbf{e}_{g_j}, & \text{nếu } (t_2 - t_1) \in [\Delta t_{\text{min}}, \Delta t_{\text{max}}] \\ -\infty, & \text{ngược lại} \end{cases}
\end{equation}
Quy tắc topo này loại trừ hơn $95\% - 98\%$ các cặp dương tính giả về mặt thị giác.

\subsection{Ước Tính Thời Gian Hành Trình và Vận Tốc Trung Bình}
Khi ứng viên đạt độ tương đồng Cosine $\ge \tau_{\text{sim}}$, hệ thống suy luận trực tiếp các chỉ tiêu giao thông hành lang:
\begin{equation}
    T_{\text{travel}} = t_2 - t_1 \quad (\text{giây}), \qquad v_{\text{journey}} = \frac{d(c_1, c_2)}{t_2 - t_1} \times 3600 \quad (\text{km/h})
\end{equation}
Giải pháp cung cấp năng lực giám sát tốc độ và thời gian hành trình liên nút giao mà không cần đầu tư hệ thống nhận diện biển số (ANPR) đắt đỏ.

\section{HƯỚNG 7: HIỂU CẢNH GIAO THÔNG TỪ VỰNG MỞ VỚI ĐỀ XUẤT VÙNG KHÔNG PHỤ THUỘC LỚP}

\subsection{Đề Xuất Vùng Vật Lý (Delta Region Proposals)}
Các mô hình nhận diện thông thường bị giới hạn trong tập danh mục khép kín cố định. Nhằm nhận diện linh hoạt các phương tiện đặc thù đô thị (xe cứu thương, xe rác, xe ba gác, chướng ngại vật rơi vãi), động cơ DeltaProposalEngine sinh ra các hộp đề xuất vùng không phụ thuộc lớp (Class-Agnostic Proposals) từ bản đồ $\Delta$ đa ngưỡng và áp dụng Non-Maximum Suppression (NMS).

\subsection{Căn Chỉnh Không Gian Thị Giác -- Ngôn Ngữ DINO-CLIP}
Mỗi vùng đề xuất $r$ được trích xuất đặc trưng qua DINO ViT và ánh xạ qua bộ chiếu đa tầng sang không gian văn bản CLIP (512 chiều):
\begin{equation}
    \mathbf{z}_r = \text{Normalize}\big(\text{Projector}(\text{DINO}(r))\big) \in \mathbb{R}^{512}
\end{equation}

Xác suất phân loại không gian mở của vùng $r$ đối với danh mục văn bản tự nhiên $c$ (được mã hóa bởi CLIP Text Encoder thành vector $\mathbf{w}_c$) được xác định bởi:
\begin{equation}
    P(c | r) = \frac{\exp\big((\mathbf{z}_r \cdot \mathbf{w}_c) / \tau\big)}{\sum_{k=1}^C \exp\big((\mathbf{z}_r \cdot \mathbf{w}_k) / \tau\big)}
\end{equation}
Hệ thống cho phép người vận hành truy vấn bất kỳ thực thể giao thông nào bằng câu mô tả văn bản tự nhiên tiếng Việt hoặc tiếng Anh mà không cần huấn luyện lại mạng.

\section{HƯỚNG 8: ƯỚC LƯỢNG ĐA THUỘC TÍNH TÌNH TRẠNG MẶT ĐƯỜNG TỰ GIÁM SÁT}

\subsection{Khai Thác Chuỗi Ảnh Nền 24 Giờ}
Kho ảnh nền tĩnh $\{I_{\text{bg}}^{h=0}, \dots, I_{\text{bg}}^{h=23}\}$ trên 608 camera ghi nhận đầy đủ sự biến đổi môi trường của hạ tầng đường bộ đô thị theo thời gian:
\begin{itemize}
    \item Điều kiện thời tiết: Mặt đường khô ráo vs. mặt đường ẩm ướt / đọng nước mưa (thể hiện qua các vùng phản chiếu gương Specular Reflection).
    \item Chu kỳ chiếu sáng: Ban ngày, hoàng hôn chạng vạng và chiếu sáng nhân tạo ban đêm.
    \item Chất lượng hạ tầng: Độ gồ ghề, ổ gà, bong tróc nứt nẻ mặt đường (phản ánh qua độ nhám gradient kết cấu).
\end{itemize}

\subsection{Kiến Trúc Đa Nhiệm và Mất Mát Mỏ Neo Vật Lý}
Mô hình RoadConditionClassifier sử dụng DINO ViT trích xuất vector kết cấu mặt đường $\mathbf{z}_{\text{road}}$ từ $[CLS]$ token và phân nhánh thành ba đầu dự đoán đồng thời:
\begin{enumerate}
    \item Chỉ số ẩm ướt mặt đường $\hat{y}_{\text{wet}} \in [0, 1]$ qua hàm kích hoạt Sigmoid.
    \item Phân loại 3 trạng thái chiếu sáng $\hat{\mathbf{y}}_{\text{illum}} \in \mathbb{R}^3$ (Đêm, Chạng vạng, Ngày).
    \item Chỉ số suy giảm kết cấu mặt đường $\hat{y}_{\text{deg}} \in [0, 1]$ phản ánh độ nứt nẻ ổ gà.
\end{enumerate}

Hàm mất mát kết hợp mỏ neo vật lý sơ cấp:
\begin{equation}
    \mathcal{L}_{\text{surface}} = \mathcal{L}_{\text{CE}}(\hat{\mathbf{y}}_{\text{illum}}, y_{\text{illum}}) + \lambda_1 \|\hat{y}_{\text{wet}} - s_{\text{specular}}\|_2^2 + \lambda_2 \|\hat{y}_{\text{deg}} - s_{\text{roughness}}\|_2^2
\end{equation}
Hệ thống tự động gom cụm phân tích PCA 2D và xuất báo cáo tình trạng toàn đô thị phục vụ công tác bảo trì hạ tầng giao thông thông minh.

\section{TỔNG HỢP VÀ ĐÁNH GIÁ CÁC HƯỚNG TIẾP CẬN}

\begin{table}[htbp]
\centering
\small
\caption{Tổng hợp phương pháp luận và tiềm năng công bố của 7 hướng nghiên cứu DINO Suite}
\label{tab:comparison}
\begin{tabular}{p{1.6cm} p{3.2cm} p{4.6cm} c p{2.8cm}}
\toprule
\textbf{Hướng} & \textbf{Tên Nghiên Cứu} & \textbf{Cơ Chế Cốt Lõi} & \textbf{Độ Mới} & \textbf{Venue Đề Xuất} \\
\midrule
\textbf{Hướng 1} & BG-Guided DINO Continual SSL & Foreground-Aware Masking (FAM) định hướng ViT che và học biểu diễn phương tiện & 4/5 & IEEE T-ITS, EAAI \\
\addlinespace
\textbf{Hướng 2} & Scene Decomposition Network & Mô hình hóa Alpha Compositing với mỏ neo giám sát nền tĩnh (Road Inpainting) & 5/5 & CVPR, ECCV, NeurIPS \\
\addlinespace
\textbf{Hướng 3} & Foreground-Enhanced Counting & Mở rộng Patch Embedding 4 kênh (RGB+$\Delta$) kết hợp Warm-Start & 3.5/5 & EAAI Journal, ITSC \\
\addlinespace
\textbf{Hướng 5} & Spatio-Temporal Density \& HCM LoS & Bi-GRU kết hợp DINO và $\Delta$-CNN với mỏ neo $\rho_{\text{phys}}$ dự đoán mật độ và LoS & 4.5/5 & IEEE T-ITS, CVPR \\
\addlinespace
\textbf{Hướng 6} & Corridor-Based Vehicle Re-ID & $\Delta$-RoI không cần detector, GeM+BNNeck, Ràng buộc không-thời gian đo tốc độ $km/h$ & 4.5/5 & IEEE T-ITS, TRB \\
\addlinespace
\textbf{Hướng 7} & Open-Vocabulary Scene Understanding & $\Delta$ proposals không phụ thuộc lớp kết hợp căn chỉnh DINO sang CLIP text & 4/5 & ECCV, WACV \\
\addlinespace
\textbf{Hướng 8} & Road Surface Condition Estimation & Đánh giá đa thuộc tính mặt đường (đọng nước, chiếu sáng, hư hại) từ ảnh nền 24h & 3.5/5 & IEEE T-ITS, TRB \\
\bottomrule
\end{tabular}
\end{table}

\section{KẾT LUẬN}
Báo cáo đã hệ thống hóa cơ sở lý thuyết về mô hình học tự giám sát DINO và hoàn thiện bảy hướng nghiên cứu có tính bổ trợ chặt chẽ nhằm khai thác cặp tín hiệu ảnh hiện trường và ảnh nền tĩnh trong giám sát giao thông đô thị. Hệ sinh thái trải rộng từ tiền huấn luyện thích ứng miền (Hướng 1), phân rã bối cảnh tự thân (Hướng 2), ước lượng lưu lượng phương tiện ít mẫu (Hướng 3), dự báo mật độ không-thời gian và cấp độ dịch vụ giao thông (Hướng 5), nhận dạng lại phương tiện theo hành lang và đo tốc độ di chuyển (Hướng 6), hiểu cảnh từ vựng mở (Hướng 7), đến đánh giá điều kiện mặt đường đô thị (Hướng 8). Việc kết hợp có hệ thống giữa các đặc tính quang học vật lý và mô hình nền tảng thị giác giúp giải quyết triệt để bài toán thiếu hụt dữ liệu gán nhãn, nâng cao độ bền vững và mở rộng khả năng suy luận thông minh cho các hệ thống giao thông đô thị quy mô lớn.

\end{document}
